\documentclass[10pt]{article}
\usepackage[margin=0.67in]{geometry}
\usepackage[utf8]{inputenc}
\usepackage[T1]{fontenc}
\usepackage[spanish]{babel}
\usepackage{xcolor}
\usepackage{booktabs}
\usepackage{tabularx}
\usepackage{hyperref}

\definecolor{navy}{HTML}{17324D}
\definecolor{teal}{HTML}{147D82}
\definecolor{coral}{HTML}{D9654F}
\definecolor{mist}{HTML}{EEF3F5}
\hypersetup{colorlinks=true,urlcolor=teal,linkcolor=navy}
\setlength{\parskip}{4pt}
\setlength{\parindent}{0pt}

\begin{document}

{\color{navy}\LARGE\textbf{Una Lima, dos realidades}}\\[-2pt]
{\large Brechas en la infraestructura urbana observable entre San Isidro y San Juan de Lurigancho}

\vspace{4pt}
\small
\textbf{Equipo:} DINAMITA --- Jeffry Arturo Hilario Quintana, Jeffry Vargas y Fatima Margarita Villon Zarate.\\
\textbf{Repositorio:} \url{https://github.com/jeffHQ/MapillaryVisualization}

\hrule\vspace{7pt}

\section*{1. Problema, contexto y propósito}

Una misma ciudad puede ofrecer experiencias urbanas muy distintas. Las personas caminan, cruzan vías, esperan transporte y se orientan mediante una combinación de veredas, cruces, señales, marcas viales e iluminación. Sin embargo, esas condiciones no siempre se presentan de la misma forma en todos los distritos de Lima Metropolitana.

El proyecto explora \textbf{cómo cambia la evidencia visual de infraestructura peatonal y vial} entre zonas urbanas comparables de San Isidro y San Juan de Lurigancho. Elegimos el contraste por su relevancia socioeconómica, no para etiquetar a sus habitantes ni para concluir que un distrito es ``mejor'' que otro. El Mapa de Pobreza del INEI brinda un contexto histórico verificable: para 2013 registra un intervalo de pobreza total de 0,0--0,3\% en San Isidro y grupos territoriales de 10,0--41,1\% en San Juan de Lurigancho.\footnote{INEI, \href{https://www.inei.gob.pe/media/DocumentosPublicos/pobreza/2013/Mapa-de-pobreza-provincial-y-distrital-2013.pdf}{Mapa de Pobreza Provincial y Distrital 2013}, cuadro 2.}

La visualización estará dirigida a estudiantes e investigadores de urbanismo, organizaciones de movilidad activa y equipos municipales. Su propósito es permitirles \textbf{explorar contrastes espaciales y formular prioridades de revisión}, no sustituir una auditoría técnica ni emitir una calificación de seguridad vial.

\textbf{¿Por qué visualizarlo?} El fenómeno es simultáneamente espacial, categórico y temporal: miles de imágenes, sus ubicaciones, sus fechas y los elementos detectados no se entienden en una tabla. Un mapa interactivo permite comparar ambas zonas, filtrar categorías y abrir la imagen que sostiene cada observación. Esto vuelve trazable la discusión: el usuario puede pasar de un patrón agregado a la evidencia visual concreta.

\section*{2. Dataset y metodología propuesta}

El dataset principal será reconstruido para esta propuesta. Se recolectarán imágenes georreferenciadas mediante la \textit{Mapillary Graph API} en \textbf{tres zonas de estudio equivalentes por distrito}, usando la misma cuadrícula, tamaño de celda, límite de imágenes y reglas de deduplicación. Trabajar con zonas comparables, en vez de descargar por completo ambos distritos, mantiene el alcance realizable y evita que el tamaño desigual de San Juan de Lurigancho domine la comparación.

\begin{tabularx}{\linewidth}{@{}p{0.24\linewidth}X@{}}
\toprule
\textbf{Fuente} & \textbf{Uso en el proyecto} \\
\midrule
Mapillary Graph API & Coordenadas, fecha de captura, secuencia, orientación, calidad, enlace de miniatura y detecciones de objetos. Mapillary es una fuente de evidencia, no el objeto de evaluación. \\
INEI & Indicador socioeconómico histórico para justificar el contraste territorial; no se usa para inferir la condición de una persona o cuadra particular. \\
Límites oficiales & Delimitación y asignación de cada imagen a su distrito y zona de estudio. \\
OpenStreetMap (si se incorpora) & Red vial como denominador complementario para comparar por tramo, no solo por cantidad de imágenes. \\
\bottomrule
\end{tabularx}

La tabla principal tendrá una fila por imagen; una tabla relacionada agregará las detecciones por \texttt{id\_imagen} y categoría. Las categorías se agruparán en: infraestructura peatonal (vereda, cruce cebra, rebaje), regulación vial (señales, semáforos, marcas), soporte urbano (postes e iluminación), movilidad y entorno. Antes del análisis se eliminarán IDs duplicados, se validará la geometría, se normalizarán fechas y categorías, y se registrará la cobertura temporal.

\textbf{Límite central.} Una detección automática en una imagen es evidencia de que un elemento fue observado desde ese punto de vista; su ausencia no demuestra que el elemento no exista ni mide su estado físico. Por ello, los indicadores se expresarán por 100 imágenes o por tramo cuando corresponda, y cada resultado podrá inspeccionarse mediante su imagen de origen.

\section*{3. Preguntas que guiarán las visualizaciones}

\begin{enumerate}
  \item ¿Cómo varía la disponibilidad y actualidad de evidencia visual entre las zonas de ambos distritos?
  \item ¿Qué categorías de infraestructura peatonal y vial aparecen con mayor frecuencia por cada 100 imágenes en cada zona?
  \item ¿Dónde se concentran espacialmente las detecciones de cruces, veredas, señalización, iluminación y marcas viales?
  \item ¿Qué contrastes persisten al comparar zonas equivalentes, en lugar de conteos brutos por distrito?
  \item ¿Qué imágenes y secuencias permiten contextualizar los patrones agregados hallados en el mapa?
  \item ¿Qué sesgos de cobertura, fecha, cámara o detección limitan la comparación y qué áreas requieren verificación posterior?
\end{enumerate}

\section*{4. Factibilidad, plan y riesgos}

El equipo reutilizará el proceso de adquisición ya desarrollado, pero con una cuadrícula simétrica y un diccionario adaptado al nuevo objetivo. Jeffry Arturo Hilario Quintana liderará la adquisición, limpieza y trazabilidad; Jeffry Vargas la preparación espacial, comparación y visualización; y Fatima Margarita Villon Zarate la narrativa, experiencia de usuario y evaluación. El mayor riesgo es la cobertura desigual de Mapillary o respuestas inestables de la API. Se mitigará con subconsultas, reintentos, un registro de adquisición, reglas de exclusión explícitas y normalización por imágenes/tramos. La entrega posterior incluirá una interfaz D3 con mapa, filtros por categoría y fecha, comparación entre zonas e inspección de evidencia.

\end{document}
